\documentclass[10pt,a4paper]{amsart}
\usepackage[margin=19mm]{geometry}
\usepackage{amsmath,amssymb,mathtools}
\usepackage[scr=rsfs,cal=euler]{mathalfa}
\usepackage{xfrac}
\usepackage[unicode,hidelinks,bookmarks=false]{hyperref}
\newcommand{\ie}{i.e.\ }
\newcommand{\cf}{cf.\ }
\newcommand{\resp}{respectively\ }
\newcommand{\Subsection}{\subsection}
\providecommand{\nobreakdash}{\nobreak\hbox{-}}
\setlength{\emergencystretch}{2em}
\multlinegap0pt
\usepackage{amssymb}
\usepackage[shortlabels]{enumitem}
\usepackage{xcolor}
\usepackage{algorithm}\usepackage{algpseudocode}
\newtheorem{definition}{Definition}
\newtheorem{prop}{Proposition}
\title{A Set-theoretic Approach to Regularity of 3D Decaying Turbulence}
\author{Min Chul Lee}
\date{Source: arXiv:2212.14734v10 (CC BY 4.0)}
\begin{document}
\maketitle

\noindent\emph{Source and licence.} A Set-theoretic Approach to Regularity of 3D Decaying Turbulence. Version arXiv:2212.14734v10 (CC BY 4.0), distributed by arXiv under the Creative Commons Attribution 4.0 International licence, http://creativecommons.org/licenses/by/4.0/, as recorded in the arXiv licence metadata for this exact version and shown on the abstract page https://arxiv.org/abs/2212.14734v10. Full licence notice: \texttt{licenses/arxiv-2212.14734-CC-BY-4.0.txt}.\par\smallskip

\noindent\textit{Editorial scope.} Counted unit: the article's prop prop (source0.tex lines 397--399). The article's complete original proof of it is delivered with it (source0.tex lines 400--411, one \texttt{proof} environment of 12 lines). No mathematics is written, repaired or re-derived: the statement and the proof are the article's own lines, in the article's own notation, and no retained line is substituted or re-flowed.  Retained uncounted as the source-local context the proof needs: lines 377--393.  Nothing was omitted from any retained span, so the package carries no omission record.  No retained line contains a citation, so the wrapper carries no bibliography and no bibliography entry.  No retained line contains a figure environment, an image-inclusion command or drawing code, so no image is delivered and the package declares no asset dependency.  Editorial formatting only, in addition: an AMS \texttt{amsart} preamble that restates the article's own declarations and macros used by the retained lines, namely the environments \texttt{definition}, \texttt{prop} on separate counters, together with the standard packages those lines need.  The retained environments print in this document as Definition 1, Proposition 1 in order of appearance, whereas the article prints the same items with the numbering of its own full document.  The complete native source of the article as distributed in the arXiv e-print for this exact version is bundled unchanged as \texttt{sources/135531/source0.tex}.
\begin{definition}
\label{partial order def}
  Let $u$ be a weak solution and $I \subset (0,\infty)$ be an open interval. For $s,t \in I$, we write $s \preceq t$ if
\begin{enumerate}
    \item $s=t$ or

    \item $s< t$ and  there exist bounded open intervals $(s_1, s_2)$ and $(t_1, t_2)$ in $I$ such that \begin{itemize}
    \item $s_2 < t_1$
      \item $s \in (s_1, s_2)$ and $t \in (t_1, t_2)$
      \item $\lVert u(t') \rVert \leq \lVert u(s') \rVert$ for almost all $s' \in (s_1, s_2)$, for almost all $t' \in (t_1, t_2)$.
  \end{itemize} 
  Note that this definition is independent of any version of $u$.
\end{enumerate}
  
   
 
\end{definition}

\begin{prop}
    The relation $\preceq$ is a partial ordering on $I$.
\end{prop}

\begin{proof}
    Reflexivity and antisymmetry are obvious from the definition. For transitivity, let $s \preceq t$ and $t \preceq w$. We only need to consider the case $s < t < w$. 

   Let $(s_1, s_2)$, $(t_1, t_2)$ and $(w_1, w_2)$ be intervals as in Definition~\ref{partial order def}. Denote by $W \subset (w_1, w_2)$ the set of $w' \in (w_1, w_2)$ satisfying $\lVert u(w') \rVert \leq \lVert u(t') \rVert$ for almost all $t' \in (t_1,t_2)$. For each $w' \in W$, let $T_{w'} \subset (t_1,t_2)$ be the set of all such $t'$. Similarly, we introduce the notations $T \subset (t_1,t_2)$ and $S_{t'} \subset (s_1,s_2)$ for each $t' \in T$.


Consider 
\begin{equation*}
\underset{t' \in T_{w'} \cap T}{\cup}  S_{t'}    
\end{equation*}
for each $w' \in W$. Since $(t_1,t_2) - \bigl( T_{w'} \cap T \bigr)$ is a null set, $T_{w'} \cap T$ is nonempty. Any $s' \in \cup_{t' \in T_{w'} \cap T}  S_{t'}$ belongs to $S_{t'}$ for some $t' \in T_{w'} \cap T$ so that $\lVert u(t') \rVert \leq \lVert u(s') \rVert$. Moreover,  $\lVert u(w') \rVert \leq \lVert u(t') \rVert$ by definition of $T_{w'}$, so that we finally obtain  $\lVert u(w') \rVert \leq \lVert u(s') \rVert$. $(s_1,s_2) - \cup_{t' \in T_{w'} \cap T}  S_{t'}$ is clearly a null set.
\end{proof}

\end{document}
